
Podobno kot smo v poglavju o curry howardovem izomorfizmupofjleleswfewkljdkljskljdfskljdfsklj
TODO: Uvod



\begin{itemize}
    \item{Pravilo aksioma} 
    \begin{prooftree}
        \AXC{}
        \RL{Ax}
        \UIC{\(x : A \vdash x : A\)}
    \end{prooftree}

    \item{Hkratni produktni tip}
    \begin{prooftree}
        \AXC{\(\G \vdash N : A\)}
        \AXC{\(\D \vdash M : B\)}
        \RL{\(\otimes I\)}
        \BIC{\(\G, \D \vdash \langle N, M \rangle : A \otimes B \)}
    \end{prooftree}
    %
    \begin{prooftree}
        \AXC{\(\G \vdash M : A \otimes B\)}
        \AXC{\(\D, x : A, y : B \vdash N : C\)}
        \RL{\(\otimes E\)}
        \BIC{\(\G, \D \vdash N : C\)}
    \end{prooftree}

    \item{Linearno funkcijski tip}
    \begin{prooftree}
        \AXC{\(\G x : A \vdash N : B \)}
        \RL{\(\multimap I\)}
        \UIC{\(\G \vdash \lambda x. N : A \multimap B\)}
    \end{prooftree}
    %
    \begin{prooftree}
        \AXC{\(\G \vdash M : A \multimap B\)}
        \AXC{\(\D \vdash N : A\)}
        \RL{\(\multimap E\)}
        \BIC{\(\G, \D \vdash M\ N\)}
    \end{prooftree}


    \item{Vsotni tip}
    \begin{center}
        \begin{bprooftree}
            \AXC{\(\G \vdash N : A\)}
            \RL{\(\oplus I_L\)}
            \UIC{\(\G \vdash \inl N :  A \oplus B\)}
        \end{bprooftree}
        \begin{bprooftree}
            \AXC{\(\G \vdash N : B\)}
            \RL{\(\oplus I_D\)}
            \UIC{\(\G \vdash \inr N :  A \oplus B\)}
        \end{bprooftree}
    \end{center}
    %
    \begin{prooftree}
        \AXC{\(\G \vdash M : A \oplus B\)}
        \AXC{\(\D_1, x: A \vdash N_1\)}
        \AXC{\(\D_2, y : B \vdash N_2\)}
        \RL{\(\oplus E\)}
        \TIC{\(\G, \D \vdash \cd{match } M \cd{ with } \inl x \mapsto N_1 | \inr y \mapsto N_2\)}
    \end{prooftree}

    \item{Alternativni produktni tip}
    \begin{prooftree}
        \AXC{\(\G \vdash N : A\)}
        \AXC{\(\G \vdash M : B\)}
        \RL{\(\& I\)}
        \BIC{\(\G \vdash N \with M : A \with B\)}
    \end{prooftree}
    %
    \begin{center}
        \begin{bprooftree}
            \AXC{\(N : A \with B\)}
            \RL{\(\& E_1\)}
            \UIC{\(\fst N : A\)}
        \end{bprooftree}
        \begin{bprooftree}
            \AXC{\(N : A \with B\)}
            \RL{\(\& E_2\)}
            \UIC{\(\snd N : B\)}
        \end{bprooftree}
    \end{center}

\end{itemize}


